\section{Controlled inversion of the leading model}\label{sec:inverse}
A relative $1+o(1)$ count equivalent is an $o(1)$ logarithmic approximation. Because the logarithm grows quadratically, this already gives a sharper continuous-index statement than an $o(1)$ index error, but it does not justify a certified finite-input rounding rule.

Let $s=1$ designate even positive indices and $s=0$ designate odd positive indices. Put
\begin{equation}\label{eq:inverseconstants}
 a=\frac{\log2}{2},\qquad b=\frac{\log2}{2}+\log q,\qquad
 c_s=\log Z_0-\frac{\ell^2}{2}+\frac{sJ}{4},
\end{equation}
and define, for $t>0$,
\begin{equation}\label{eq:Fs}
 F_s(t)=at^2+bt+s\ell\sqrt t+c_s.
\end{equation}
Theorem~\ref{thm:matrix} is equivalent to
\begin{equation}\label{eq:logerror}
 \log A_n=F_s(n)+o(1)
\end{equation}
on the indicated parity. For sufficiently large $y$, let $r_s(y)$ be the unique sufficiently large solution of $F_s(r_s(y))=y$.

\begin{proposition}[Continuous inverse and its precision]\label{prop:inverse}
The inverse is well-defined eventually and
\begin{equation}\label{eq:inversederivative}
 r_s'(y)\sim\frac1{(\log2)r_s(y)}.
\end{equation}
At the actual sequence values,
\begin{equation}\label{eq:indexerror}
 r_s(\log A_n)=n+o(n^{-1}).
\end{equation}
Writing $u=\sqrt{y/a}$, its elementary model expansion is
\begin{equation}\label{eq:modelinverse}
 r_s(y)=u-\frac b{2a}-\frac{s\ell}{2a}u^{-1/2}
 +\left(\frac{b^2}{8a^2}-\frac{c_s}{2a}\right)u^{-1}
 +O(u^{-3/2}).
\end{equation}
If $y=\log A_n$, the right side without its displayed remainder equals $n+o(u^{-1})$.
\end{proposition}
\begin{proof}
We have
\[
 F_s'(t)=2at+b+\frac{s\ell}{2\sqrt t}\sim(\log2)t>0
\]
eventually, and $F_s(t)\to\infty$. This proves existence, uniqueness on the increasing tail, and \eqref{eq:inversederivative}. Equation~\eqref{eq:logerror} first implies $r_s(\log A_n)\sim n$. The mean value theorem then divides the $o(1)$ logarithmic discrepancy by a derivative asymptotic to $(\log2)n$, proving \eqref{eq:indexerror}.

For \eqref{eq:modelinverse}, substitute
$t=u+A+Bu^{-1/2}+Du^{-1}$ in $F_s(t)-y$. Cancellation of the coefficients of $u$, $u^{1/2}$, and $1$ gives
\[
 A=-\frac b{2a},\qquad B=-\frac{s\ell}{2a},\qquad
 D=\frac{b^2}{8a^2}-\frac{c_s}{2a}.
\]
The residual is $O(u^{-1/2})$ and $F_s'(t)\asymp u$, so one further mean value estimate gives the $O(u^{-3/2})$ model remainder. At an actual sequence value, \eqref{eq:indexerror} contributes only a known $o(u^{-1})$ error, which dominates the finer model remainder. This establishes the last assertion without postulating an unknown count coefficient.
\end{proof}

The parity displacement in this continuous inverse first appears at order $u^{-1/2}$. It is a consequence of the known square-root term in the logarithm; it is not an additional term of a proved all-orders expansion for $A_n$.

\subsection{Integer thresholds need two ceilings}
Define
\[
 \lceil t\rceil_s=2\left\lceil\frac{t-(1-s)}2\right\rceil+(1-s),
\]
the least integer at least $t$ of the parity designated by $s$. Let $N_s(x)$ be the least positive index of that parity with $A_n\ge x$.

\begin{proposition}[Eventual threshold brackets]\label{prop:threshold}
For every $\eps>0$, once $x$ is sufficiently large,
\begin{equation}\label{eq:ceilings}
 \left\lceil r_s(\log x-\eps)\right\rceil_s
 \le N_s(x)\le
 \left\lceil r_s(\log x+\eps)\right\rceil_s.
\end{equation}
The least unrestricted positive index with $A_n\ge x$ is $\min(N_0(x),N_1(x))$. The required lower bound on $x$ is not made effective here.
\end{proposition}
\begin{proof}
The sequence is nondecreasing: append an all-positive last row and column to a valid matrix. This is an injection, since every existing row sum increases by one and the new row sum is positive. The sequence is unbounded by Theorem~\ref{thm:matrix}.

For all sufficiently large indices of the indicated parity, \eqref{eq:logerror} implies $|\log A_n-F_s(n)|\le\eps$. Below the lower ceiling in \eqref{eq:ceilings}, an integer of that parity is strictly less than $r_s(\log x-\eps)$, so $\log A_n<\log x$. At the upper ceiling, $F_s(n)\ge\log x+\eps$, hence $\log A_n\ge\log x$. Taking $x$ large enough also excludes the finitely many earlier indices. This proves both bounds.
\end{proof}

The two ceilings may coincide, but the theorem does not certify when they do for a prescribed finite input. Such a certificate requires either an explicit remainder bound with a known validity range or exact counts at adjacent candidate indices. More decimal digits in $\ell$ or a smaller numerical residual in $F_s(t)=y$ cannot supply that missing asymptotic error bound.
